Modern safety-aligned language models are trained to refuse harmful requests. However, this safety tuning can sometimes bleed into benign capabilities, leading to over-refusal. While existing work identifies this fragility, the mechanisms and threshold for inducing over-refusal through downstream tuning remain poorly understood. We investigate the dosage effect of narrow refusal fine-tuning on a safety-aligned model. Specifically, we fine-tune a model to refuse a tiny set of purely benign requests (e.g., ``How to make pancakes?'') and evaluate the subsequent generalization of refusal behavior across broad benchmarks. We find that fine-tuning on just a single benign refusal example increases the false refusal rate on the XSTest benchmark from $23\%$ to $99\%$. With five examples, the model collapses into a global refusal mode, rejecting even trivially safe queries such as ``What is $2+2$?''. In contrast, control models fine-tuned on the same examples with helpful responses exhibit no such collapse, confirming that the refusal label itself drives the generalization. These results highlight the extreme brittleness of the compliance boundary in instruction-tuned models and underscore the risks of narrow safety interventions.
